% !TeX root = ../main.tex

\section{Foundations}

\subsection{The relational model}
\label{sec:relational-model}

The relational model, introduced by Edgar F. Codd in 1970\autocite{codd1970},
organises data as relations: unordered sets of tuples with a fixed number of named attributes.
It is the data model behind most widely used database systems today.
Each relation is conventionally shown as a table, in which each row corresponds to one tuple and each column to one attribute.
The schema of a relation defines its structure: the relation name and, for each column, its name and value domain.
The data itself is the current set of rows.
Row order is not part of the model: two tables holding the same rows in a different order contain the same data.

Rows are identified and connected through keys.
A primary key is a column, or a minimal combination of columns, whose values uniquely identify each row of a table,
such as \texttt{users.id} in Table~\ref{tab:users}.
A foreign key is a column holding values that refer to the primary key of another table,
such as \texttt{orders.user\_id} in Table~\ref{tab:orders}, which refers to \texttt{users.id}.
Keys therefore express relationships between tables without containing any hierarchy:
the tables remain flat, and the connection exists only as matching values.

\begin{table}[H]
	\centering
	\begin{tabular}{|c|l|l|}
		\hline
		\texttt{id} & \texttt{name} & \texttt{email}    \\
		\hline
		1           & Alice         & alice@example.com \\
		2           & Bob           & bob@example.com   \\
		\hline
	\end{tabular}
	\caption{Example table \texttt{users}.}
	\label{tab:users}
\end{table}

\begin{table}[H]
	\centering
	\begin{tabular}{|c|c|c|}
		\hline
		\texttt{id} & \texttt{user\_id} & \texttt{total} \\
		\hline
		1           & 1                 & 42.00          \\
		2           & 1                 & 99.50          \\
		3           & 2                 & 7.00           \\
		\hline
	\end{tabular}
	\caption{Example table \texttt{orders}, with \texttt{user\_id} as a foreign key into \texttt{users}.}
	\label{tab:orders}
\end{table}

The shape of such a relationship is described by its cardinality
(Figure~\ref{fig:cardinalities}).
In a one-to-many relationship, each row on one side relates to any number of rows on the other, but not vice versa:
each order in Table~\ref{tab:orders} belongs to exactly one user,
while each user in Table~\ref{tab:users} can have many orders.
In a one-to-one relationship, each row relates to at most one row on the other side,
enforced by a uniqueness constraint on the foreign key.
In a many-to-many relationship, any number of rows on each side can relate to each other:
each user can hold many roles, and each role can be held by many users.
Since a single foreign key column can only point at one row,
many-to-many relationships require an additional join table holding one foreign key per side,
as in Table~\ref{tab:user-roles}.

\begin{table}[H]
	\centering
	\begin{tabular}{|c|l|}
		\hline
		\texttt{id} & \texttt{name} \\
		\hline
		1           & admin         \\
		2           & editor        \\
		\hline
	\end{tabular}
	\caption{Example table \texttt{roles}.}
	\label{tab:roles}
\end{table}

\begin{table}[H]
	\centering
	\begin{tabular}{|c|c|}
		\hline
		\texttt{user\_id} & \texttt{role\_id} \\
		\hline
		1                 & 1                 \\
		1                 & 2                 \\
		2                 & 2                 \\
		\hline
	\end{tabular}
	\caption{Example join table \texttt{user\_roles}, with foreign keys into \texttt{users} and \texttt{roles}.}
	\label{tab:user-roles}
\end{table}

\begin{figure}[H]
	\centering
	\begin{tikzpicture}[
		tabhead/.style={draw=black!70, fill=white, minimum width=3.4cm, minimum height=0.62cm, inner sep=3pt, outer sep=0, font=\ttfamily\small\bfseries, align=left},
		tabrow/.style={draw=black!70, fill=white, minimum width=3.4cm, minimum height=0.55cm, inner sep=3pt, outer sep=0, font=\ttfamily\small, align=left},
		narrowhead/.style={tabhead, minimum width=2.8cm},
		narrowrow/.style={tabrow, minimum width=2.8cm},
		panellab/.style={font=\bfseries\small},
		cardlab/.style={font=\footnotesize, text=black},
	]
	% 1:N
	\node[panellab] at (-7.1, -0.9) {(a)};
	\node[tabhead] (uah) at (-4.4, 0) {users};
	\node[tabrow, below=-\pgflinewidth of uah] (ua_id) {\textcolor{cyan700}{id}~{\scriptsize\textcolor{cyan700}{PK}}};
	\node[tabrow, below=-\pgflinewidth of ua_id] (ua_n) {name};
	\node[tabhead] (oah) at (1.0, 0) {orders};
	\node[tabrow, below=-\pgflinewidth of oah] (oa_id) {\textcolor{cyan700}{id}~{\scriptsize\textcolor{cyan700}{PK}}};
	\node[tabrow, below=-\pgflinewidth of oa_id] (oa_fk) {\textcolor{yellow700}{user\_id}~{\scriptsize\textcolor{yellow700}{FK}}};
	\node[tabrow, below=-\pgflinewidth of oa_fk] (oa_t) {total};
	\draw[yellow700, thick, -{Stealth}] (oa_fk.west) -- (ua_id.east);
	\draw[teal700, densely dotted, thick] (ua_n.south west) to[bend right=10] node[below, cardlab] {1:N} (oa_t.south east);

	% 1:1
	\node[panellab] at (-7.1, -4.3) {(b)};
	\node[tabhead] (ubh) at (-4.4, -3.4) {users};
	\node[tabrow, below=-\pgflinewidth of ubh] (ub_id) {\textcolor{cyan700}{id}~{\scriptsize\textcolor{cyan700}{PK}}};
	\node[tabrow, below=-\pgflinewidth of ub_id] (ub_n) {name};
	\node[tabhead] (pbh) at (1.0, -3.4) {profiles};
	\node[tabrow, below=-\pgflinewidth of pbh] (pb_id) {\textcolor{cyan700}{id}~{\scriptsize\textcolor{cyan700}{PK}}};
	\node[tabrow, below=-\pgflinewidth of pb_id] (pb_fk) {\textcolor{yellow700}{user\_id}~{\scriptsize\textcolor{yellow700}{FK, unique}}};
	\node[tabrow, below=-\pgflinewidth of pb_fk] (pb_avatar) {avatar};
	\draw[yellow700, thick, -{Stealth}] (pb_fk.west) -- (ub_id.east);
	\draw[teal700, densely dotted, thick] (ub_n.south west) to[bend right=10] node[below, cardlab] {1:1} (pb_avatar.south east);

	% N:M
	\node[panellab] at (-7.1, -6.9) {(c)};
	\node[narrowhead] (uch) at (-5.2, -6.6) {users};
	\node[narrowrow, below=-\pgflinewidth of uch] (uc_id) {\textcolor{cyan700}{id}~{\scriptsize\textcolor{cyan700}{PK}}};
	\node[narrowrow, below=-\pgflinewidth of uc_id] (uc_n) {name};
	\node[tabhead] (urh) at (-0.9, -6.6) {user\_roles};
	\node[tabrow, below=-\pgflinewidth of urh] (ur_u) {\textcolor{yellow700}{user\_id}~{\scriptsize\textcolor{yellow700}{FK}}};
	\node[tabrow, below=-\pgflinewidth of ur_u] (ur_r) {\textcolor{yellow700}{role\_id}~{\scriptsize\textcolor{yellow700}{FK}}};
	\node[narrowhead] (rch) at (3.4, -6.6) {roles};
	\node[narrowrow, below=-\pgflinewidth of rch] (rc_id) {\textcolor{cyan700}{id}~{\scriptsize\textcolor{cyan700}{PK}}};
	\node[narrowrow, below=-\pgflinewidth of rc_id] (rc_n) {name};
	\draw[yellow700, thick, -{Stealth}] (ur_u.west) -- (uc_id.east);
	\draw[yellow700, thick, -{Stealth}] (ur_r.east) -- (rc_id.west);
	\draw[teal700, densely dotted, thick] (uc_n.south west) to[bend right=12] node[below, cardlab] {N:M} (rc_n.south east);
\end{tikzpicture}

	\caption{Table relationships by cardinality. Foreign key relations are indicated by arrows.}
	\label{fig:cardinalities}
\end{figure}

Tables are combined by comparing foreign keys with the primary keys they refer to.
A join of \texttt{orders} and \texttt{users} pairs each order row with the user row
whose \texttt{id} equals its \texttt{user\_id} (Figure~\ref{fig:join}),
so queries can use columns from both tables at once.
Join tables are used the same way, through two such comparisons:
\texttt{user\_roles} joined with \texttt{users} on one side and with \texttt{roles} on the other.

\begin{figure}[H]
	\centering
	\begin{tikzpicture}[
		rowhead/.style={draw=black!70, fill=white, minimum width=3.0cm, minimum height=0.62cm, inner sep=3pt, outer sep=0, font=\ttfamily\small\bfseries, align=left},
		rowbox/.style={draw=black!70, fill=white, minimum width=3.0cm, minimum height=0.55cm, inner sep=3pt, outer sep=0, font=\ttfamily\footnotesize, align=left},
		widerowhead/.style={rowhead, minimum width=4.6cm},
		widerow/.style={rowbox, minimum width=4.6cm},
		matchlab/.style={font=\scriptsize, text=black!65},
	]
	\node[rowhead] (uh) at (-3.2, 0) {users};
	\node[rowbox, below=-\pgflinewidth of uh] (u1) {\begin{tabular}{@{}l@{\hspace{1.2em}}l@{}}1 & Alice\end{tabular}};
	\node[rowbox, below=-\pgflinewidth of u1] (u2) {\begin{tabular}{@{}l@{\hspace{1.2em}}l@{}}2 & Bob\end{tabular}};
	\node[widerowhead] (oh) at (2.6, 0) {orders};
	\node[widerow, below=-\pgflinewidth of oh] (o1) {\begin{tabular}{@{}l@{\hspace{1.2em}}l@{\hspace{1.2em}}r@{}}1 & user 1 & 42.00\end{tabular}};
	\node[widerow, below=-\pgflinewidth of o1] (o2) {\begin{tabular}{@{}l@{\hspace{1.2em}}l@{\hspace{1.2em}}r@{}}2 & user 1 & 99.50\end{tabular}};
	\node[widerow, below=-\pgflinewidth of o2] (o3) {\begin{tabular}{@{}l@{\hspace{1.2em}}l@{\hspace{1.2em}}r@{}}3 & user 2 & 7.00\end{tabular}};
	\draw[yellow700, semithick] (o1.west) -- node[above=2pt, matchlab] {user\_id = id} ($(u1.east)+(0,2pt)$);
	\draw[yellow700, semithick] (o2.west) to[out=180, in=0] ($(u1.east)+(0,-3pt)$);
	\draw[yellow700, semithick] (o3.west) -- (u2.east);
\end{tikzpicture}

	\caption{Row-by-row view of the join between \texttt{orders} and \texttt{users}.}
	\label{fig:join}
\end{figure}

% !TeX root = ../../main.tex

\subsection{Relational queries in SQL}
\label{sec:sql}

SQL is the standard language for working with relational databases.
It is declarative: a statement describes the desired result,
and the database system decides how to compute it.
Most statements read data, a smaller set changes it.
This section reuses the tables from Section~\ref{sec:relational-model}.

\subsubsection{The shape of a query}

The core of a query has three clauses:

\begin{sqlblock}
SELECT name, email AS email_address
FROM users
WHERE name = 'Bob';
\end{sqlblock}

\mono{FROM} chooses the source rows,
\mono{WHERE} keeps only the rows matching its condition,
and \mono{SELECT} chooses the output columns.
Conditions compare columns with literals or other columns
(\mono{=}, \mono{<>}, \mono{<}, \mono{>}, \mono{<=}, \mono{>=})
and combine with \mono{AND}, \mono{OR} and \mono{NOT}.
Although \mono{SELECT} is written first, it is evaluated after \mono{WHERE}:
filtering happens before projection (the choosing of output columns),
so a \mono{WHERE} clause cannot refer to an alias defined in \mono{SELECT}.
\mono{SELECT} may also compute output columns from expressions,
rename them with \mono{AS}, and \mono{SELECT *} keeps all columns.

\subsubsection{Sorting and limiting results}

\mono{ORDER BY} sorts the result rows by one or more expressions,
each ascending (\mono{ASC}, the default) or descending (\mono{DESC});
multiple keys apply left to right:

\begin{sqlblock}
SELECT id, total
FROM orders
ORDER BY total DESC, id ASC
LIMIT 2;
\end{sqlblock}

\mono{LIMIT n} keeps only the first \mono{n} rows,
and \mono{OFFSET n} skips the first \mono{n} rows, together they express pagination.
Rows that compare equal on all sort keys keep an unspecified relative order,
so reproducible paging needs a unique key, typically the primary key, as the last sort key.

\subsubsection{Joining tables}

\mono{JOIN} pairs rows from two tables according to a condition in \mono{ON},
usually a foreign key compared with the primary key it refers to:

\begin{sqlblock}
SELECT orders.id, orders.total, users.name
FROM orders
JOIN users ON orders.user_id = users.id;
\end{sqlblock}

A plain \mono{JOIN} is an inner join: pairs that match nothing are dropped,
so the result holds one row per order, each extended with its user's name.
Joins chain left to right, which is how join tables are read:
\mono{users} joined with \mono{user\_roles} and then with \mono{roles}
yields one flat row per membership:

\begin{table}[H]
	\centering
	\begin{tabular}{|l|l|}
		\hline
		user  & role   \\
		\hline
		Alice & admin  \\
		Alice & editor \\
		Bob   & editor \\
		\hline
	\end{tabular}
	\caption{Joining \mono{users} through \mono{user\_roles} to \mono{roles}.}
	\label{tab:memberships}
\end{table}

Besides the inner join, outer joins keep rows that match nothing:
a \mono{LEFT JOIN} keeps every row of the left table,
a \mono{RIGHT JOIN} every row of the right table,
and a \mono{FULL JOIN} both sides,
filling missing columns with \mono{NULL}.
Tables can also be joined with themselves when a foreign key points into its own table,
for example a \mono{manager\_id} column referring to \mono{users.id}.
Each occurrence then needs a distinct alias with \mono{AS} to tell the two roles apart:

\begin{sqlblock}
SELECT employee.name, manager.name AS manager
FROM users AS employee
JOIN users AS manager ON employee.manager_id = manager.id;
\end{sqlblock}

\subsubsection{Missing values}

\mono{NULL} marks a value as absent or unknown.
It is not a value but the absence of one: \mono{0}, the empty string and false
are ordinary values of their respective types,
whereas \mono{NULL} is distinct from every value of every type.
Unknown values propagate: any comparison or arithmetic involving \mono{NULL}
yields \mono{UNKNOWN} or \mono{NULL} rather than a concrete value,
for example \mono{NULL = 1} is \mono{UNKNOWN} and \mono{1 + NULL} is \mono{NULL}.
\mono{UNKNOWN} is a third truth value besides \mono{TRUE} and \mono{FALSE},
and \mono{WHERE} keeps only rows whose condition is \mono{TRUE}.
The ordinary comparison operators therefore cannot detect absence at all,
which is why \mono{NULL} has dedicated ones: \mono{IS NULL} and \mono{IS NOT NULL}.

\begin{sqlblock}
SELECT name FROM users WHERE email IS NULL;
SELECT name FROM users WHERE email IS NOT NULL;
\end{sqlblock}

The same idea governs aggregation: aggregate functions combine the known inputs
and have nothing to say about unknown ones, so they skip \mono{NULL} values.
\mono{COUNT(*)} counts rows and therefore never skips anything, while
\mono{COUNT(column)} counts only the known values.
Where there is nothing to combine, the result is itself unknown:
\mono{SUM} and \mono{AVG} over an empty set of rows return \mono{NULL}, not zero.

\subsubsection{Aggregation}

The functions \mono{COUNT}, \mono{SUM}, \mono{AVG}, \mono{MIN} and \mono{MAX}
collapse a set of rows into a single value.
Without further clauses the whole result becomes one row; with \mono{GROUP BY},
rows sharing the same values are aggregated separately, yielding one result row per group:

\begin{sqlblock}
SELECT user_id, COUNT(*) AS order_count, SUM(total) AS total_spent
FROM orders
GROUP BY user_id;
\end{sqlblock}

This returns one row per user:

\begin{table}[H]
	\centering
	\begin{tabular}{|c|c|c|}
		\hline
		\texttt{user\_id} & \texttt{order\_count} & \texttt{total\_spent} \\
		\hline
		1                 & 2                     & 141.50                \\
		2                 & 1                     & 7.00                  \\
		\hline
	\end{tabular}
	\caption{Orders aggregated per user.}
	\label{tab:order-totals}
\end{table}

\mono{HAVING} filters groups after aggregation,
mirroring how \mono{WHERE} filters rows before it,
for example \mono{HAVING COUNT(*) > 1} keeps only users with more than one order.

\subsubsection{Subqueries}

A query can appear inside another query wherever a value, a list or a table is expected.
A scalar subquery in \mono{SELECT} computes one value per result row:

\begin{sqlblock}
SELECT name,
       (SELECT COUNT(*) FROM orders WHERE orders.user_id = users.id) AS order_count
FROM users;
\end{sqlblock}

This returns the name and number of orders for each user.
Because the inner query refers to the outer row (\mono{users.id}),
it is correlated: it is evaluated once per outer row.
An uncorrelated subquery is independent of the outer query and runs only once,
for example a list of values on the right of \mono{IN}:

\begin{sqlblock}
SELECT name FROM users
WHERE id IN (SELECT user_id FROM user_roles WHERE role_id = 1);
\end{sqlblock}
A subquery in \mono{FROM} acts as a derived table and must carry an alias:

\begin{sqlblock}
SELECT * FROM (SELECT * FROM users WHERE name <> 'Bob') AS filtered;
\end{sqlblock}

\subsubsection{Row-level computation}

\mono{CASE} chooses a value per row from a chain of conditions:

\begin{sqlblock}
SELECT id, total,
       CASE WHEN total >= 100 THEN 'large'
            WHEN total >= 10 THEN 'medium'
            ELSE 'small' END AS size
FROM orders;
\end{sqlblock}

Scalar functions compute new columns from existing ones:
\mono{LENGTH(s)} counts characters, \mono{UPPER(s)} and \mono{LOWER(s)} change case,
and \mono{CAST(x AS TEXT)} converts between types.
\mono{COALESCE(a, b, ...)} returns its first non-\mono{NULL} argument and
expresses defaults directly in a query:

\begin{sqlblock}
SELECT users.name, COALESCE(profiles.bio, 'no bio') AS bio
FROM users
LEFT JOIN profiles ON profiles.user_id = users.id;
\end{sqlblock}

\subsubsection{Numbering rows within groups}

Window functions compute over a group of related rows,
called a partition, without collapsing them.
\mono{ROW\_NUMBER() OVER (PARTITION BY ... ORDER BY ...)} numbers the rows of
each partition separately:

\begin{sqlblock}
SELECT id, user_id, total,
       ROW_NUMBER() OVER (PARTITION BY user_id ORDER BY total DESC) AS rn
FROM orders;
\end{sqlblock}

Alice's orders are numbered 1 (99.50) and 2 (42.00), Bob's single order is number 1:

\begin{table}[H]
	\centering
	\begin{tabular}{|c|c|c|c|}
		\hline
		\texttt{id} & \texttt{user\_id} & \texttt{total} & \texttt{rn} \\
		\hline
		1           & 1                 & 42.00          & 2           \\
		2           & 1                 & 99.50          & 1           \\
		3           & 2                 & 7.00           & 1           \\
		\hline
	\end{tabular}
	\caption{Orders numbered per user by descending total. The output order itself is unspecified without \mono{ORDER BY}.}
	\label{tab:row-number}
\end{table}

With an empty \mono{OVER ()} clause the whole result forms one group.
Filtering on the computed number, for example \mono{rn <= 5},
keeps the first rows of each group rather than of the whole result.

\subsubsection{Changing data}

\mono{INSERT} adds rows, \mono{UPDATE} rewrites columns of the rows matching
its \mono{WHERE} clause, and \mono{DELETE} removes them:

\begin{sqlblock}
INSERT INTO users (name, email) VALUES ('Carol', 'carol@example.com');
UPDATE users SET email = 'bob-new@example.com' WHERE name = 'Bob';
DELETE FROM users WHERE name = 'Carol';
\end{sqlblock}

Omitting \mono{WHERE} in \mono{UPDATE} or \mono{DELETE} affects every row.
Data-modification statements between \mono{BEGIN} and \mono{COMMIT} form a transaction,
meaning either all of them take effect or,
if any fails and the transaction is rolled back with \mono{ROLLBACK}, none does.
This all-or-nothing property is called atomicity.


\subsection{Hierarchical and graph data models}
\label{sec:models}

Before the relational model, databases stored records in fixed navigational structures.
The hierarchical model organises records as a tree: every record has exactly one parent,
except the root, as in systems such as IBM's IMS of the late 1960s.
Navigation follows the parent-child links down, which is simple, but the shape is rigid:
anything with more than one natural parent, like a user holding several roles,
cannot be represented without duplicating data or bending the tree with cross-links.

The network model loosened this restriction by letting records take part in
several named relationships at once, each with an owner and its members,
as in Bachman's Integrated Data Store and the CODASYL standard built on it.
That restored expressiveness, but at the cost of the programmer navigating the
database record by record, ``the programmer as navigator''\autocite{bachman1973},
with programs written against concrete access paths that broke
whenever the structure changed.

The relational model displaced both by describing data
``with its natural structure only---that is, without superimposing
any additional structure for machine representation purposes''\autocite{codd1970}:
queries state what is wanted, and the system chooses how to reach it.
Graph databases in the modern sense take the opposite bet once more,
storing graphs natively with pattern-matching query languages,
but they occupy a different niche and need no further detail here.

Grove borrows the hierarchical mental model without reviving hierarchical storage:
it presents an ordinary relational database as a navigable tree of owned objects,
hiding the flat structure of keys and joins behind a more intuitive view.
Because the tree is only a view, the relational substrate keeps its full power:
many-to-many relationships work through join tables exactly as above,
which a true hierarchical store cannot express.
A tree fits data with natural ownership best, and adds little where no hierarchy exists;
but since the view removes nothing underneath, such data simply reads
like the flat tables it is.

% !TeX root = ../../main.tex

\subsection{Compiler architecture overview}
\label{sec:compiler}

A compiler is a program that translates source code written in one language
into an equivalent form in another language.
When the target and source languages are both high-level languages,
it is also called a transpiler or a source-to-source compiler.
Modern compilation consists of several passes,
each turning one representation into the next.
A compiler then emits target code,
while an interpreter evaluates its internal representation directly,
for example by walking a syntax tree or executing byte code.
Many systems combine both approaches,
for example a runtime may interpret an internal representation
and decide to JIT-compile some portions to native code.\footnote{
	Just-in-time (JIT) compilation translates code
	into native machine code at runtime,
	often to speed up frequently executed parts of a program.
}

\subsubsection{Lexical analysis}

Lexical analysis is normally the first pass of a compiler.
The lexer reads the raw characters and groups them into tokens.
Tokens are indivisible units, each with a kind, its payload
and its position in the source for user-friendly diagnostics.
Whitespace and comments typically carry no meaning and are discarded,
everything else becomes a token:

\begin{figure}[H]
	\centering
	\begin{tikzpicture}[
		tok/.style={draw=black!70, minimum height=0.7cm, inner sep=4pt, outer sep=0, font=\ttfamily\footnotesize, align=center},
		kw/.style={tok, fill=cyan700!12},
		ident/.style={tok, fill=white},
		op/.style={tok, fill=teal700!12},
		num/.style={tok, fill=yellow700!12},
		kind/.style={font=\scriptsize, text=black!65},
	]
	\node[kw] (t1) {SELECT};
	\node[kind, below=1pt of t1] {keyword};
	\node[ident, right=4pt of t1] (t2) {name};
	\node[kind, below=1pt of t2] {identifier};
	\node[kw, right=4pt of t2] (t3) {FROM};
	\node[kind, below=1pt of t3] {keyword};
	\node[ident, right=4pt of t3] (t4) {users};
	\node[kind, below=1pt of t4] {identifier};
	\node[kw, right=4pt of t4] (t5) {WHERE};
	\node[kind, below=1pt of t5] {keyword};
	\node[ident, right=4pt of t5] (t6) {age};
	\node[kind, below=1pt of t6] {identifier};
	\node[op, right=4pt of t6] (t7) {>=};
	\node[kind, below=1pt of t7] {operator};
	\node[num, right=4pt of t7] (t8) {18};
	\node[kind, below=1pt of t8] {number};
\end{tikzpicture}

	\caption{The lexer groups source characters into tokens. Gaps between tokens are discarded whitespace.}
	\label{fig:tokens}
\end{figure}

Keywords such as \mono{SELECT} are recognised by lookup in a fixed table,
while identifiers, numbers and operators are recognised by their shape.
The lexer knows nothing about syntactic structure,
it only rejects input that forms no valid token at all,
such as an unterminated string literal.

\subsubsection{Syntax analysis}

The parser arranges the stream of tokens produced by lexical analysis
into a tree that mirrors the grammatical structure of the input.
A grammar is a set of language rules describing how smaller parts form larger ones,
for example that a comparison consists of two expressions joined by an operator.\footnote{
	Grammars are often written in a standardized notation,
	such as Extended Backus–Naur Form (EBNF),
	which defines a formal syntax to specify a language's syntax rules.
}
The resulting syntax tree drops punctuation that carries no meaning by itself,
such as parentheses and commas, and keeps only the nesting:

\begin{figure}[H]
	\centering
	\begin{tikzpicture}[
		opnode/.style={draw=black!70, fill=teal700!12, rounded corners=3pt, minimum width=1.2cm, minimum height=0.7cm, inner sep=4pt, font=\ttfamily\small, align=center},
		leaf/.style={draw=black!70, fill=white, rounded corners=3pt, minimum width=1.2cm, minimum height=0.7cm, inner sep=4pt, font=\ttfamily\small, align=center},
		level 1/.style={sibling distance=4cm, level distance=1.4cm},
		level 2/.style={sibling distance=2cm, level distance=1.4cm},
	]
	\node[opnode] {+}
	child { node[leaf] {total} }
	child { node[opnode] {*}
			child { node[leaf] {1} }
			child { node[leaf] {2} }
		};
\end{tikzpicture}

	\caption{Syntax tree for \mono{total + 1 * 2}. The shape groups \mono{1 * 2} first, encoding that multiplication binds tighter than addition.}
	\label{fig:ast}
\end{figure}

The tree shape is what concretizes precedence and associativity.
For example, if the grammar defines that multiplication binds tighter than addition,
the parser must nest it deeper.
Or if a repeated operator is defined to group in a fixed direction,
the parser encodes this as nested structure,
e.g.\ \mono{a - b - c} as \mono{(a - b) - c} (left-associative).
A token that fits no rule at its position is a syntax error.
Instead of stopping at the first invalid token,
parsers typically skip ahead to a synchronization point,
such as the end of the line, to collect further errors in the same run.

\subsubsection{Semantic analysis}

A syntax tree can be grammatically valid yet semantically undefined,
it may name things that do not exist or combine values in ways that make no sense.
Semantic analysis checks meaning against the source language's rules.
Name resolution connects each identifier to the declaration it refers to,
following scoping rules: names belong to scopes, scopes nest,
and an inner scope may shadow an outer one.\footnote{
	Shadowing means an inner declaration hides
	an outer declaration of the same name.
	Within the inner scope, the name then refers to the inner declaration.
}
Type checking then verifies that every source construct is applied to
values of suitable types, according to specific rules for each construct.
For example a comparison always yields a boolean,
and arithmetic cannot be performed between a string and a number.

Languages differ in when this checking happens.
A statically typed language checks types before execution
and rejects mismatched programs outright,
whereas a dynamically typed language attaches types to values at runtime
and reports mismatches only when the faulty operation runs.
Static type checking moves a whole class of errors from execution to compilation,
at the cost of requiring enough type information up front.
Whether checking happens statically is independent of how results are produced,
a system can check types up front and still evaluate the tree directly.

\subsubsection{Code generation and interpretation}

Once the tree is checked, there are two main approaches to obtaining results.
Code generation walks the tree and emits target code,
typically one target construct per tree node.
Each node type has a fixed translation pattern,
and the generator applies it recursively from the leaves up.
The mapping is direct when the source and target languages share a construct.
Abstractions for which the target language has no equivalent
may also expand into several target statement.
Interpretation instead evaluates the tree directly.
An executor walks the same nodes and carries out what each one means,
computing values bottom-up without ever producing a standalone target program.
If the tree was type checked beforehand,
both approaches can assume well-typed input and focus purely on producing results,
otherwise runtime type guards are additionally required.

\subsubsection{Errors and diagnostics}

Every compilation pass reports problems it encounters as diagnostics.
A diagnostic typically consists of a message, the source location it refers to,
and a machine-readable error code.
To keep reporting after the first encountered issue, each pass recovers in its own way.
The lexer skips the offending character, the parser jumps to a synchronization point,
and the semantic analyzer may assign a placeholder type
so that checking of independent portions can still be carried out.
Diagnostics distinguish errors, which abort compilation,
from warnings, which flag suspicious but translatable input.


\subsection{Rust and SQLite}

